\documentclass[a4paper,10pt,notitlepage]{article}
\usepackage{ctex,geometry,graphicx,tikz,setspace,paralist,fancyhdr,caption}
\geometry{
	left=1.5cm,
	right=1.5cm,
	top=1.4cm,
	bottom=2cm,}
\newcommand{\rec}{
	\begin{tikzpicture}[remember picture,overlay]
		% 绘制边框
		\draw[line width=1.2pt] ([xshift=0.5cm,yshift=0.5cm] current page.south west) rectangle ([xshift=-0.5cm,yshift=-0.8cm] current page.north east);
	\end{tikzpicture}
	
}
\pagestyle{fancy}
\fancyhf{}
\fancyhead[C]{\rec}  % 在页眉中绘制图形
\fancyhead[L]{附录}
\fancyhead[R]{实验七 \quad 积分与微分电路}
\fancyfoot[C]{\thepage}
\begin{document}
	\large
	\onehalfspacing
	\begin{figure}[h]
		\raggedright
		\includegraphics{预习报告/1.png}
	\end{figure}
	\centering
	{\Huge\textbf{模电实验报告}\par}
	\vspace{0.2cm}
	{\huge{实验内容：积分与微分电路}\par}
	\raggedright
	\vspace{0.3cm}
	\begin{centering}
		{\large 院系：电子与信息工程学院\hfill 学号：22309080\hfill 审批：\hspace{2cm} \par
			专业：通信工程\hfill 实验人：梁倍铭\hfill 日期：2023年11月28日 \par}
	\end{centering}
	\vspace{0.3cm}
	\section*{一、实验目的}
	\begin{enumerate}
		\item 学会用运算放大器组成积分微分电路。
		\item 学会积分微分电路的特点及性能。
	\end{enumerate}
	\section*{二、原理简介}
	\begin{enumerate}
		\item 积分电路\par 
		实验电路如图7-1 所示
		\begin{figure}[h]
			\centering
			\includegraphics[width=0.5\textwidth]{1.png}
			\caption*{图7-1 积分电路}
		\end{figure}
		反积分电路：
		$$U_o=-\frac{1}{R_1C}\int_{t_0}^{t}U_i(t)\mathrm{d}t+U_o(t_0)$$
		积分电路输出电压是输入电压的积分，随着不同的输入电压，输出电压也表现为不同
		的形式。电路除了进行积分运算外，很多情况下应用在波形变换电路中。
		\item 微分电路\par 
		实验电路如图7-2 所示。
		\begin{figure}[h]
			\centering
			\includegraphics[width=0.5\textwidth]{2.png}
			\caption*{图7-2 微分电路}
		\end{figure}
		微分电路电路理想分析得到公式：
		$$U_o(t)=-RC\frac{\mathrm{d}U_i(t)}{\mathrm{d}t}$$
		输出电压是输入电压的微分。
		\item 积分—微分电路\par 
		实验电路如图7-3 所示。
		\begin{figure}[h]
			\centering
			\includegraphics[width=0.5\textwidth]{3.png}
			\caption*{图7-3 积分--微分电路}
		\end{figure}
		先积分后微分可以对输入信号一个大致的还原。
	\end{enumerate}
	\section*{三、实验步骤和内容}
	\begin{enumerate}
		\item 积分电路\par 
		按图7-1 连接电路。
		\begin{enumerate}
			\item 第一个实验课上老师说直接跳过
			\item 使图7-1 中积分电容改为0.1μ，断开K1，Vi 分别输入100Hz 幅值为2V 方
			波正弦波信号，观察Vi 和Vo 大小及相位关系，并记录波形。波形图如下
			\begin{figure}[h]
				\centering
				\begin{minipage}{0.3\textwidth}
					\centering
					\includegraphics[width=\textwidth]{4.png}
					\caption*{图7-4 仿真波形}
				\end{minipage}
				\qquad
				\begin{minipage}{0.3\textwidth}
					\centering
					\includegraphics[width=\textwidth]{1.jpg}
					\caption*{图7-5 实验波形}
				\end{minipage}
				\qquad
				\begin{minipage}{0.3\textwidth}
					\centering
					\includegraphics[width=\textwidth]{2.jpg}
					\caption*{图7-6 实验失真波形}
				\end{minipage}
			\end{figure}
			\item 改变输入的频率，观察Vi 与Vo 的相位、幅值关系。关系变化如表7-1
			\begin{table}[h]
				\centering
				\begin{tabular}{|c|c|c|c|c|}
					\hline
					频率(Hz) & 100 & 200 & 300 & 400\\
					\hline
					仿真幅值 & 3.15 & 1.565 & 1.055 & 0.79\\
					\hline 
					实验幅值 & 3.06 & 1.542 & 1.057 & 0.82\\
					\hline
				\end{tabular}
				\caption*{表7-1 频率幅值表}
			\end{table}
		\end{enumerate}
		\item 微分电路\par
		\begin{enumerate}
			\item 输入有效值为lV，f =160Hz 三角波（正弦波）信号，用示波器观察Vi 与Vo
			波形并测量输出电压。
			\begin{figure}[h]
				\centering
				\begin{minipage}{0.3\textwidth}
					\centering
					\includegraphics[width=\textwidth]{5.png}
					\caption*{图7-7 仿真波形}
				\end{minipage}
				\qquad
				\begin{minipage}{0.3\textwidth}
					\centering
					\includegraphics[width=\textwidth]{3.jpg}
					\caption*{图7-8 实验波形}
				\end{minipage}
			\end{figure}
			\item 改变三角波（正弦波）频率(20HZ～400HZ)，观察Vi 与Vo 的相位、幅值变化
			情况并记录。幅值变化关系如表7-2，相位都是相差$\frac{\pi}{2}$
			\begin{table}[h]
				\centering
				\begin{tabular}{|c|c|c|c|c|c|c|c|}
					\hline
					频率(Hz) & 20 & 50 & 100 & 150 & 200 & 300 & 400 \\
					\hline
					仿真幅值 & 0.277 & 0.7 & 1.38 & 2.08 & 2.775 & 4.16 & 5.55\\
					\hline 
					实验幅值 & 0.258 & 0.68 & 1.372 & 2.04 & 2.658 & 4.15 & 5.440 \\
					\hline
				\end{tabular}
				\caption*{表7-2 频率幅值表}
			\end{table}
			\item 输入V= ±5V，f =200Hz 的方波信号，用示波器观察Vo 波形，按上述步骤重
			复实验。
			\begin{figure}[h]
				\centering
				\begin{minipage}{0.3\textwidth}
					\centering
					\includegraphics[width=\textwidth]{6.png}
					\caption*{图7-9 仿真波形}
				\end{minipage}
				\qquad
				\begin{minipage}{0.3\textwidth}
					\centering
					\includegraphics[width=\textwidth]{4.jpg}
					\caption*{图7-10 实验波形}
				\end{minipage}
			\end{figure}
		\end{enumerate}
		\item 积分--微分电路\par 
		按图7-3 连接电路。
		\begin{enumerate}
			\item 在Vi 输入f =200Hz，V=±6V 的方波信号，用示波器观察Vi 和V0 的波形并
			记录。观察到的波形如下
			\begin{figure}[h]
				\centering
				\begin{minipage}{0.3\textwidth}
					\centering
					\includegraphics[width=\textwidth]{7.png}
					\caption*{图7-11 仿真波形}
				\end{minipage}
				\qquad
				\begin{minipage}{0.3\textwidth}
					\centering
					\includegraphics[width=\textwidth]{5.jpg}
					\caption*{图7-12 实验波形}
				\end{minipage}
			\end{figure}
			
			\item 将f 改为500Hz 重复上述实验。
			\begin{figure}[h]
				\centering
				\begin{minipage}{0.3\textwidth}
					\centering
					\includegraphics[width=\textwidth]{8.png}
					\caption*{图7-13 仿真波形}
				\end{minipage}
				\qquad
				\begin{minipage}{0.3\textwidth}
					\centering
					\includegraphics[width=\textwidth]{6.jpg}
					\caption*{图7-14 实验波形}
				\end{minipage}
			\end{figure}
		\end{enumerate}
	\end{enumerate}
	\section*{四、实验器材}
	1、 实验箱 2、数字万用表 3、函数信号发生器 4、交流毫伏表
	5、双踪示波器
	\section*{五、实验结果分析}
	\begin{enumerate}
		\item 对于图7-1和7-2的分析
		\begin{enumerate}
			\item 分析图7-1 电路，若输入正弦波，Vo 与Vi 相位差是$-90^{\circ}$当输入信号为100Hz 有 效
			值为2V 时，Vo=7.4v
			\item 图7-2 电路，若输入方波，Vo 与Vi 相位差是$90^{\circ}$当输入信号为160Hz 幅值为lV
			时，输出Vo=6v
		\end{enumerate}
		\item 积分电路的特点
		\begin{enumerate}
			\item 输出信号与输入信号的关系是积分关系，输出电压是输入电压的积分。
			\item 对于低频信号的积分，积分电路表现出较好的性能。它能够积累来自输入信号的全部变化，并消除高频成分
			\item 积分电路通常用于信号处理、滤波、计算总量、微分方程的求解等应用。
		\end{enumerate}
		\item 微分电路的特点
		\begin{enumerate}
			\item 输出信号与输入信号的关系是微分关系，输出电压是输入电压的变化率。
			\item 对于高频信号的微分，微分电路表现出较好的性能。它能够放大高频信号的微小变化并滤除低频成分。
			\item 微分电路通常用于信号处理、滤波、峰值检测等应用。
		\end{enumerate}
	\end{enumerate}
\end{document}